\itemhead{习题4.1 · 第6题（教材p.104）}{XJ-B1-C04-HW-P104-Y1259}
\label{XJ-B1-C04-HW-P104-Y1259}
{\footnotesize 湘教版教材p.104 / PDF第111页 / 习题4.1 / 第6题}
\field{原题}{化简：
\subq{1}$\left(\pi^{\sqrt2+1}\right)^{\sqrt2-1}$；
\subq{2}$\left(\left(\sqrt3\right)^{\sqrt[4]{64}}\right)^{\sqrt{25}}$。}
\field{参考答案}{（1）$\pi$；（2）$3^{5\sqrt2}$，因为 $\sqrt[4]{64}=2\sqrt2,\sqrt{25}=5$。以当前题库的根指数及括号为准。}
{\small \field{关联知识}{P02 有理数、实数指数幂与运算（教材帮p.144、145 / PDF153、154）}
\field{方法与易错点}{统一为指数幂，分开处理系数与各底数的指数；最终复查定义条件及数值精度。 易错：负指数误作负数；幂的幂指数相加；约分丢禁点；过早四舍五入。}
\field{答案依据}{答案册相关题组：PDF第31页（册内页30）；按当前题库原题逐项重算，不按旧版题号机械对套。}}
